\chapter{Регионы и~города}

Остров удобно делить на~пять ландшафтных областей:
столица \place{Лас-Пальмас}, зелёный \textit{Norte}, изрезанный
\textit{Oeste}, солнечный \textit{Sur} и~горное ядро \textit{Centro}.

\section{Лас-Пальмас-де-Гран-Канария \sightt{}}

\factline{380~000 жителей. Автобус №1, 12, 17 вдоль набережной.}

Столица --- шумный портовый город с~двумя центрами притяжения.
На~севере --- старый квартал \place{Вегета} с~жёлто-охристыми
фасадами XVI~в.: \sight{Катедраль-де-Санта-Ана}, начатая в~1497~г.
и~объединяющая позднюю готику и~канарский ренессанс;
\sight{Каса-де-Колон}, восстановленный губернаторский дом, где,
по~преданию, ночевал Колумб по~пути в~Новый Свет; и~рынок
\textit{Mercado de Vegueta} (пятница–суббота, утро).

\marginfact{Плайя-\\де-Лас-Кантерас:\\3,2~км мелкого\\песка,\\безопасный\\риф}

В~трёх километрах севернее тянется \sight{Плайя-де-Лас-Кантерас},
пожалуй, лучший городской пляж Европы: подводный риф
\textit{La~Barra} гасит океанскую волну, позволяя купаться
с~маленькими детьми. Вечером --- променад \textit{Paseo} с~сотней
кафе.

\begin{detour}
Из~порта \textit{Puerto de~la~Luz} паромы \textit{Fred.~Olsen}
за~час идут на~Тенерифе, за~два --- на~Фуэртевентуру.
\end{detour}

\section{Север: зелёные склоны}

\subsection*{\place{Терор}}

\factline{14~км от~Лас-Пальмаса. Автобус №216.}

Живописный городок с~резными деревянными балконами
(\textit{balcones canarios}) и~\sight{базиликой Вирхен-дель-Пино}
--- покровительницы острова. Воскресный рынок известен
чорисо и~сладостями местного монастыря.

\subsection*{\place{Арукас}}

Город построен из~чёрного вулканического базальта; внушительная
\textit{неоготическая} церковь Сан-Хуан-Баутиста (1909) кажется
собором крупного европейского города. Здесь же --- единственная
на~Канарах винокурня рома, \textit{Arehucas} (XIX~в.).

\subsection*{\place{Агаэте}}

Рыбацкое городок у~западного входа в~долину \textit{Valle de
Agaete} --- единственного места в~Европе, где в~промышленных
объёмах выращивают кофе. Отсюда же отходит паром на~Тенерифе
(30~мин).

\section{Запад: отвесные скалы}

Западное побережье --- самое труднодоступное. Дорога \textit{GC-200}
цепляется за~скалы над~океаном; перед \place{Пуэрто-де-ла-Алдеа}
она прорезает пять тоннелей. Севернее, под~массивом
\place{Тамадаба}, --- нетронутый сосновый лес и~лучшие смотровые
площадки острова (\textit{Mirador del Balcón}).

\section{Юг: песок и солнце}

\subsection*{\place{Маспаломас} и~\place{Плайя-дель-Инглес}}

Курортная агломерация, обязанная своим рождением архитектурному
конкурсу 1961~г. Главная природная достопримечательность ---
\sightt{Дюны Маспаломаса}: 404~га подвижного песка, охраняемая
территория. Тропа вдоль пляжа от~маяка \textit{Faro de~Maspalomas}
до~\textit{Playa del Inglés} --- 6~км.

\subsection*{\place{Пуэрто-де-Моган}}

Иногда называют \textit{«канарской Венецией»} --- белые домики
с~синими ставнями стоят вдоль узких каналов. Здесь заканчиваются
многие западные тропы и~отходят лодки на~\textit{Puerto Rico}
и~для наблюдения за~дельфинами.

\section{Центр: вершина острова}

\subsection*{\place{Техеда}}

\factline{Высота 1050~м. Автобус №305 от~Лас-Пальмаса, 1\,ч~30\,мин.}

Белоснежная деревня в~чаше древней кальдеры; Мигель де~Унамуно
назвал здешний пейзаж «окаменевшей бурей».
Со~смотровой \sight{Cruz de~Tejeda} (1580~м) открывается классический
канарский вид: \sightt{Роке-Нубло}, \sight{Роке-Бентайга}
и,~при~ясной погоде, вулкан \textit{Тейде} на~Тенерифе.

\subsection*{\place{Артенара}}

Самая высоко расположенная деревня острова (1270~м). Жители по~сей
день обитают в~\textit{casas cueva} --- пещерных домах, вырубленных
в~вулканическом туфе; часть из~них превращена в~небольшие
гостиницы.
